Die Position (4) ist die kapitalisierte jährliche Ausgabe von 900 Thalern, aus der der variable Teil der Witwenpension zu bestreiten ist. Diese Summe wird unter allen berechtigten Witwen zu gleichen Teilen verteilt, wenn ihre Anzahl 18 oder mehr beträgt; ist die Anzahl kleiner, so erhält jede 50 Thaler. Daraus geht hervor, dass die Kasse im letzteren Fall (der auch jetzt vorliegt) eine Ersparnis erzielt, die in der Bilanz nicht mitgerechnet ist und für eine etwas längere Zeit auch gar nicht im Voraus berechnet werden kann. Jedenfalls ist dies aber nur ein vorübergehender Vorteil, der später, wenn sich die Folgen der jetzigen großen Ausdehnung der Genossenschaft erst entwickelt haben werden, selten oder vielleicht niemals wieder vorkommen wird.

Als Durchschnittswerte der Nebenausgaben habe ich angenommen
\begin{center}
\small\setlength{\tabcolsep}{2pt}
\begin{tabular}{r r l r l r l p{196pt}}
&28&Thl.&4&Ggr&—&Pf.&für Prozesskosten, soweit nicht erstattet,\\
&61&,,&8&,,&9&,,&für Baukosten\\
&133&,,&9&,,&5&,,&für Rechnungsführung und Gebühren für Abschriften\\
&136&,,&6&,,&3&,,&für Verluste\\
Zusammen&359&Thl.&4&Ggr.&5&Pf.&
\end{tabular}
\end{center}
\noindent\hspace*{\parindent}Diese Ansätze beruhen auf den Erfahrungen der letzten 21 oder 20 Jahre. Der Rechnung von 1845 konnten nur Erfahrungen von 15 oder 14 Jahren zugrunde gelegt werden, die die Gesamtsumme 323 Thl. 13 Ggr. ergeben hatten. Die erste Zahl erbringt kapitalisiert 10262 Thl. 10 Ggr. Dieser Betrag bildet zusammen mit 2644 Thl. 4 Ggr. (als mittlerem Wert des nicht ertragbringenden Vermögensteils nach einem 17-jährigen Durchschnitt) die Position (5). Für die zweite Zahl war übrigens in der Rechnung von 1845 nach einem 11-jährigen Durchschnitt der nahezu gleiche Wert 2635 Thl. 18 Ggr. 4 Pf. angenommen worden.

Die Position (6) entsteht durch die Kapitalisierung der jährlichen Einnahme aus dem Pachtzins der Universitätsapotheke (1000 Thaler).
